% Annexe de la note « Fusionner les impôts sur le revenu » : niches, inventaire des impôts sur le revenu, vérifications
% Compilation depuis fiscalite/ir_csg/ : pdflatex annexe-ir-csg.tex (deux passes)
\documentclass[11pt,a4paper]{article}
% Préambule commun au brief et à la revue de littérature sur la taxe sur la valeur foncière.
% Même composition que le cours (cours/refaire-l-impot.tex) ; figures « nues » produites par
%   FIG_NU=1 python lvt/graphiques_lvt.py
\usepackage[T1]{fontenc}
\usepackage[utf8]{inputenc}
\usepackage{lmodern}
\usepackage[french]{babel}
\usepackage[a4paper,left=2.8cm,right=2.8cm,top=2.6cm,bottom=2.8cm]{geometry}
\usepackage{graphicx}
\usepackage[export]{adjustbox}
\usepackage{booktabs}
\usepackage{tabularx}
\usepackage{array}
\usepackage{xcolor}
\usepackage{titlesec}
\usepackage[labelfont=bf,labelsep=period,font=small,singlelinecheck=false,
            justification=raggedright,skip=6pt]{caption}
\usepackage{fancyhdr}
\usepackage{enumitem}
\usepackage{microtype}
\usepackage{textcomp}
\usepackage{amsmath}
\providecommand{\euro}{\texteuro}
\usepackage[hidelinks]{hyperref}

\definecolor{encre}{HTML}{1A1A19}
\definecolor{bleu}{HTML}{003399}
\definecolor{vert}{HTML}{1F7A4D}
\definecolor{ocre}{HTML}{C79100}
\definecolor{brique}{HTML}{993333}
\definecolor{filet}{HTML}{C8C6BE}

\titleformat{\section}{\normalfont\large\bfseries\color{bleu}}{\thesection.}{0.6em}{}
\titleformat{\subsection}{\normalfont\normalsize\bfseries\color{bleu}}{\thesubsection}{0.6em}{}
\titlespacing*{\section}{0pt}{2.4ex plus 1ex minus .2ex}{1.2ex plus .2ex}
\titlespacing*{\subsection}{0pt}{2ex plus 1ex minus .2ex}{0.8ex plus .2ex}

\pagestyle{fancy}
\fancyhf{}
\renewcommand{\headrulewidth}{0.4pt}
\fancyhead[R]{\small\thepage}

\setlength{\parindent}{0pt}
\setlength{\parskip}{0.62em}
\linespread{1.05}
\renewcommand{\arraystretch}{1.2}
\setlist{itemsep=2pt,topsep=2pt,parsep=0pt,leftmargin=1.4em}

\newenvironment{encadre}[1]
  {\par\medskip\noindent\begin{minipage}{\linewidth}%
   \color{encre}\setlength{\parskip}{0.5em}%
   \hrule height 0.8pt width \linewidth \vspace{6pt}%
   {\footnotesize\bfseries\color{vert}\MakeUppercase{#1}}\par\vspace{-2pt}%
   \small}
  {\par\vspace{4pt}\hrule height 0.4pt width \linewidth\end{minipage}\par\medskip}

\newcommand{\figbloc}[3]{%
  \begin{figure}[htbp]\centering
  \caption{#2}\label{#3}
  \includegraphics[max width=\linewidth,max height=0.40\textheight]{../figures/nu/#1.pdf}
  \IfFileExists{../figures/nu/#1_note.tex}{\par\vspace{4pt}%
    \parbox{\linewidth}{\raggedright\scriptsize\color{encre}\setlength{\parskip}{2pt}%
    \input{../figures/nu/#1_note.tex}}}{}%
  \end{figure}}

\newcommand{\pc}{\,\%}

\fancyhead[L]{\small\itshape Fusionner les impôts sur le revenu : annexe}
\begin{document}
\renewcommand{\tablename}{Tableau}
\thispagestyle{empty}
{\footnotesize\color{ocre}\bfseries ANNEXE\quad·\quad NOVEMBRE 2026\par}
\vspace{6pt}
{\fontsize{18}{22}\selectfont\bfseries\color{bleu}Fusionner les impôts sur le revenu : données, niches et vérifications\par}
\vspace{5pt}
{\small Annexe de la note « Fusionner les impôts sur le revenu : une fausse bonne idée ? ». Colin Baget.\par}
\vspace{5pt}
{\color{bleu}\hrule height 1.2pt}
\vspace{8pt}

\textbf{Les niches supprimées.} Les dépenses fiscales de l'IR valent 44~Md€ (\emph{Voies et moyens}, PLF 2026).
L'assiette commune en supprime sûrement 24~Md€ : les réductions et crédits d'impôt (15,5~Md€, dont emploi à
domicile 7,0, dons 2,2, garde d'enfants 1,8, investissement locatif 1,5, outre-mer 1,3), l'abattement de
10~\% sur les pensions (5,3~Md€), les demi-parts spéciales (2,9~Md€) et l'abattement des plus de 65~ans
avec quelques déductions du revenu global (0,5~Md€, dont 0,4 pour l'abattement). Selon la définition retenue des revenus, jusqu'à 12~Md€ de plus (épargne salariale, heures
supplémentaires, régimes du foncier). Disparaissent aussi l'abattement de 10~\% des salariés, les
déductions d'épargne retraite et de pensions alimentaires, la décote et le quotient. L'allègement de 3,8~Md€
sur les revenus de source étrangère, imposé par les conventions fiscales, reste.

\textbf{Les impôts sur le revenu en France, et ce que la fusion remplace} (revenus 2024).

\begin{table}[!ht]\centering\scriptsize
\begin{tabularx}{\linewidth}{@{}>{\raggedright\arraybackslash}X>{\raggedright\arraybackslash}p{3.3cm}>{\raggedright\arraybackslash}p{3.8cm}@{}}
\toprule
Prélèvement & Rendement & Dans la fusion ? \\
\midrule
IR sur les salaires, pensions, revenus des indépendants, loyers et allocations chômage (après réductions et crédits) & 80~Md€ & Oui \\
Contribution exceptionnelle sur les hauts revenus (CEHR, 3 et 4~\%) & Moins de 2~Md€, comprise dans l'IR de la DGFiP & Oui, absorbée ; la part due sur les revenus du capital est reportée sur les autres revenus \\
Contribution différentielle sur les hauts revenus (CDHR, taux minimal de 20~\%) & 0,4~Md€ en 2025 ; créée pour les revenus 2025 & Absente de nos données 2024 \\
CSG et CRDS sur les revenus d'activité et de remplacement & 132 et 7,4~Md€ & Oui \\
IR sur les revenus du capital au prélèvement forfaitaire unique (12,8~\%) & Environ 11,5~Md€ & Non \\
CSG, CRDS et prélèvement de solidarité sur les revenus du capital (9,2, 0,5 et 7,5~\%) & Environ 33~Md€ (calcul à partir de l'assiette de la CRDS) & Non \\
CASA des retraités (0,3~\%) & Environ 1~Md€ (estimation) & Non \\
CRDS sur les prestations familiales et les jeux & 0,5~Md€ & Non \\
Cotisations sociales & --- & Non : l'assiette est le revenu après cotisations \\
\bottomrule
\end{tabularx}
\end{table}

\begin{table}[!ht]\centering\scriptsize
\begin{tabularx}{\linewidth}{@{}>{\raggedright\arraybackslash}p{3.2cm}>{\raggedright\arraybackslash}X>{\raggedright\arraybackslash}p{4.1cm}@{}}
\toprule
Ce qui est testé & Source de comparaison & Résultat \\
\midrule
Composition de la rubrique « salaires » de l'IRCOM & Totaux des cases 2042, revenus 2024 (DGFiP) & 925 Md€ par les cases contre 930 dans l'IRCOM ; pensions 391 contre 393 \\
IR actuel & DGFiP, Statistiques n°~41 (2025) & 91,7 Md€, après réductions et crédits, comme la DGFiP (92) \\
CSG sur les salaires & Commission des comptes de la Sécurité sociale, mai 2026 (privé 72,2, public 19,8~Md€) & 91,5 reconstitués contre 92,0 \\
CRDS & CADES, rapport annuel 2024 (9,2~Md€, dont 7,7 dans le champ) & 7,4~Md€ reconstitués \\
Pensions brutes & Eurostat (412~Md€ en 2023) & 406~Md€ reconstitués pour 2024 \\
Barème de l'IR (exercice 1) & OpenFisca-France, paramètres 2024 & Seuils 11~497, 29~315, 83~823 et 180~294~€ \\
Revenu après cotisations & OpenFisca-France, cas-types & Salaires : 88,7~\% du brut \\
Taux d'un salarié seul & OCDE, \emph{Taxing Wages 2026} & Écart de 0,1 à 0,4~point \\
Partage dans les couples & Insee Première 1492 & Part moyenne de la femme : 36~\% \\
Chiffres de la note & \texttt{proposition.py} les recalcule tous & Sortie \texttt{sorties/proposition.txt} \\
\bottomrule
\end{tabularx}
\end{table}


\end{document}
